\documentclass[12pt,a4paper]{amsart}

\usepackage{amsmath,amssymb,amsthm,mathtools,enumitem,booktabs,array,hyperref}
\usepackage[utf8]{inputenc}
\usepackage[T1]{fontenc}
\usepackage{geometry}
\geometry{margin=2.5cm}

% -------------------------------------------------------
% Theorem environments
% -------------------------------------------------------
\newtheorem{theorem}{Theorem}[section]
\newtheorem{lemma}[theorem]{Lemma}
\newtheorem{corollary}[theorem]{Corollary}
\newtheorem{proposition}[theorem]{Proposition}
\theoremstyle{definition}
\newtheorem{definition}[theorem]{Definition}
\newtheorem{example}[theorem]{Example}
\newtheorem{remark}[theorem]{Remark}
\newtheorem{conjecture}[theorem]{Conjecture}

% -------------------------------------------------------
% Custom commands
% -------------------------------------------------------
\newcommand{\N}{\mathbb{N}}
\newcommand{\Z}{\mathbb{Z}}
\newcommand{\Q}{\mathbb{Q}}
\newcommand{\R}{\mathbb{R}}
\newcommand{\C}{\mathbb{C}}
\newcommand{\Ac}{\mathcal{A}}
\newcommand{\Bc}{\mathcal{B}}
\newcommand{\Cc}{\mathcal{C}}
\newcommand{\At}{\widehat{\mathcal{A}}}
\newcommand{\dr}{\,\mathrm{d}}
\newcommand{\dbar}{\bar{\partial}}
\newcommand{\KZ}{\mathrm{KZ}}
\newcommand{\ord}{\mathrm{ord}}
\DeclareMathOperator{\degree}{deg}
\DeclareMathOperator{\End}{End}
\DeclareMathOperator{\tr}{tr}
\DeclareMathOperator{\STU}{STU}

% -------------------------------------------------------
\begin{document}
% -------------------------------------------------------

\title{The Kontsevich Integral and Vassiliev Invariants:\\
       Classification of Knots via Chord Diagrams}

\author{Michael Fuhrmann}
\address{Specialist Mathematics Project, Build 170}
\email{specialist-maths@localhost}
\date{\today}

\begin{abstract}
Vassiliev (finite-type) knot invariants, introduced by Vassiliev in 1990 and
reformulated combinatorially by Birman and Lin, are among the most powerful
and structured invariants in low-dimensional topology.  The Kontsevich
integral $Z(K)$, constructed by Maxim Kontsevich in 1993, is a universal
Vassiliev invariant: it is a formal power series in the graded algebra $\mathcal{A}$
of chord diagrams modulo the $4T$ and $1T$ relations, taking values such that
every Vassiliev invariant factors through~$Z$.  In this paper we develop
the full framework: chord diagrams and the $4T$/$1T$ relations; weight systems
and their relation to Lie algebras; the Kontsevich integral as an iterated
Chen iterated integral; the Drinfeld associator $\Phi_{\mathrm{KZ}}$ and its
role in regularising the integral at crossings; the coloured Jones polynomial
as a specialisation; the open conjecture that $Z$ is a complete invariant;
the Wheeling theorem (Bar-Natan--Le--Thurston); and the Le--Murakami--Ohtsuki
invariant.  All unproven statements are carefully labelled as conjectures.
\end{abstract}

\maketitle

\tableofcontents

% -------------------------------------------------------
\section{Introduction and Historical Overview}
% -------------------------------------------------------

The classification of knots --- the topological type of embeddings
$K: S^1 \hookrightarrow S^3$ up to ambient isotopy --- is a central
problem of low-dimensional topology.  Classical invariants such as the
Alexander polynomial (1928) and the Jones polynomial (1984) distinguish
many knots, but neither is a complete invariant.

The discovery of the Jones polynomial~\cite{jones1985} triggered a revolution:
Witten~\cite{witten1989} interpreted it via $3$-dimensional Chern-Simons
gauge theory with gauge group $\mathrm{SU}(2)$, suggesting that a ``complete''
knot invariant might arise from quantum field theory.  Simultaneously,
Vassiliev~\cite{vassiliev1990} introduced finite-type invariants via the
topology of discriminant spaces (spaces of singular knots).

Kontsevich~\cite{kontsevich1993} unified these approaches by constructing
the \emph{Kontsevich integral}: a single formal-power-series-valued invariant
that contains all Vassiliev invariants as coefficients of its expansion.
Kontsevich's construction draws on the theory of iterated integrals of Chen,
on the Drinfeld associator from quantum groups, and on combinatorial structures
of chord diagrams.

\subsection*{Organisation of the paper}

Section~\ref{sec:vassiliev} defines Vassiliev invariants via the
Vassiliev resolution and singular knots.
Section~\ref{sec:chord} introduces chord diagrams, the $4T$ and $1T$ relations,
and the algebra $\Ac$ of chord diagrams.
Section~\ref{sec:weight} develops weight systems and their Lie-algebraic
construction.
Section~\ref{sec:KI} constructs the Kontsevich integral via iterated integrals.
Section~\ref{sec:convergence} addresses convergence and the Drinfeld associator.
Section~\ref{sec:associator} treats the Knizhnik-Zamolodchikov associator $\Phi_{\mathrm{KZ}}$.
Section~\ref{sec:jones} relates the Kontsevich integral to the coloured Jones polynomial.
Section~\ref{sec:universality} proves the universality of the Kontsevich integral.
Section~\ref{sec:completeness} discusses the completeness conjecture.
Section~\ref{sec:barley} presents the wheeling theorem.
Section~\ref{sec:LMO} introduces the Le-Murakami-Ohtsuki invariant.
Section~\ref{sec:chern} connects the theory to Chern-Simons field theory.
Section~\ref{sec:open} lists open problems.

% -------------------------------------------------------
\section{Vassiliev (Finite-Type) Invariants}\label{sec:vassiliev}
% -------------------------------------------------------

\begin{definition}[Knot invariant]
A \emph{knot} is an ambient isotopy class of smooth embeddings $S^1 \to S^3$
(or $\R^3$).  A \emph{knot invariant} is a map $v: \{\text{knots}\} \to \mathbf{k}$
(where $\mathbf{k}$ is a commutative ring, typically $\Q$ or $\C$) invariant
under ambient isotopy.
\end{definition}

\begin{definition}[Singular knots and the Vassiliev skein relation]
A \emph{singular knot with $n$ double points} is a smooth immersion
$K: S^1 \looparrowright S^3$ with exactly $n$ transverse self-intersections.
Extend any knot invariant $v$ to singular knots by the \emph{Vassiliev
(finite-difference) relation}:
\[
  v\!\left(\,\begin{tikzpicture}[baseline=-2pt,scale=0.35]
    % crossing symbol
  \end{tikzpicture}\,\right)
  = v(K_+) - v(K_-),
\]
where $K_+$ (resp.\ $K_-$) is obtained by resolving the double point
positively (resp.\ negatively).  More precisely:
\[
  v\Bigl(\,\raisebox{-3pt}{\includegraphics[height=12pt]{crossing}}\,\Bigr)
  = v(K_+) - v(K_-).
\]
In symbolic notation, for a double point $\times$ we write
$v(\times) = v(K_+) - v(K_-)$.
\end{definition}

Since the above requires a picture, we give the formal algebraic definition:

\begin{definition}[Vassiliev skein relation, algebraic form]\label{def:vassiliev}
Let $v: \{\text{knots}\} \to \mathbf{k}$ be a knot invariant.  Extend $v$ to
singular knots inductively: for a singular knot $K_\times$ with a marked
double point,
\[
  v(K_\times) \;:=\; v(K_+) - v(K_-),
\]
where $K_+, K_-$ are the two resolutions of the double point (positive and
negative crossing).  Then $v$ is a \emph{Vassiliev invariant of type $\leq n$}
(or \emph{finite-type invariant of order $n$}) if $v$ vanishes on all singular
knots with more than $n$ double points:
\[
  v(K_\times) = 0 \quad \text{for all singular knots with } > n \text{ double points}.
\]
\end{definition}

\begin{example}[Vassiliev invariants of small order]
\begin{enumerate}[label=(\alph*)]
  \item Every invariant of type $0$ is constant (i.e.\ detects only the unknot
        from all knots if it's non-constant).
  \item Every invariant of type $1$ is constant (since there are no
        type-$1$ non-trivial invariants).
  \item The first non-trivial Vassiliev invariant has type $2$; it is the
        coefficient $a_2$ in the Conway polynomial $\nabla(K)(z) = \sum a_{2i} z^{2i}$.
        Equivalently, $a_2(K) = \frac{1}{6}(J_2(K)|_{q=1})$ where $J_2$ is
        the coloured Jones polynomial for the 2-dimensional representation.
  \item The Casson invariant $\lambda(K)$ (for knots in homology spheres) is
        a Vassiliev invariant of type $2$.
\end{enumerate}
\end{example}

\begin{theorem}[Birman-Lin 1993, Vassiliev invariants and Taylor expansion]
\label{thm:birmanlin}
Vassiliev invariants of order $\leq n$ are precisely those knot invariants
whose ``$n$-th derivative at the unknot'' is well-defined, where the
derivative is taken via the Vassiliev skein relation.  They form a
filtered vector space $\mathcal{V}_0 \subseteq \mathcal{V}_1 \subseteq
\mathcal{V}_2 \subseteq \cdots$ with $\bigcup_n \mathcal{V}_n = \mathcal{V}$.
\end{theorem}

% -------------------------------------------------------
\section{Chord Diagrams and the Algebra $\mathcal{A}$}\label{sec:chord}
% -------------------------------------------------------

\begin{definition}[Chord diagram]
A \emph{chord diagram of order $n$} is an oriented circle (representing the
domain $S^1$) with $n$ unordered pairs of distinct points connected by
\emph{chords} (arcs in the interior of the disc).  Two chord diagrams are
considered the same if they are related by an orientation-preserving
diffeomorphism of the circle.  Let $\Ac_n$ denote the $\Q$-vector space
spanned by chord diagrams of order $n$.
\end{definition}

\begin{definition}[$4T$ relation]
The \emph{four-term relation ($4T$)} in $\Ac_n$ is:
\[
  D_1 - D_2 + D_3 - D_4 = 0,
\]
where $D_1, D_2, D_3, D_4$ are four chord diagrams that agree outside
a fixed local region, and in that region differ by the four ways to connect
two strands with one additional chord (while keeping all other chords fixed).
More precisely, for two chords sharing an endpoint:
\[
  [ab] - [ac] + [bc] - [bd] = 0,
\]
in a region with strands $a, b, c, d$ from left to right.
\end{definition}

\begin{definition}[$1T$ relation (framing independence)]
The \emph{one-term relation ($1T$)} says that any chord diagram containing
an \emph{isolated chord} (a chord whose two endpoints are adjacent on the
circle, with no other chord endpoints between them) equals zero.
\end{definition}

\begin{definition}[Chord diagram algebra $\mathcal{A}$]
The \emph{chord diagram algebra} is the graded $\Q$-algebra:
\[
  \Ac = \bigoplus_{n=0}^\infty \Ac_n, \quad
  \Ac_n = \frac{\Q[\text{chord diagrams of order } n]}{4T + 1T \text{ relations}}.
\]
The multiplication is given by connected sum of chord diagrams:
$D \cdot D' = D \# D'$ (concatenation along the circle).
\end{definition}

\begin{theorem}[Dimensions of $\mathcal{A}_n$]
The first few dimensions are:
\[
  \dim \Ac_0 = 1, \quad \dim \Ac_1 = 0, \quad \dim \Ac_2 = 1,
  \quad \dim \Ac_3 = 1, \quad \dim \Ac_4 = 3, \quad \dim \Ac_5 = 4.
\]
The generating function $\sum_n (\dim \Ac_n) t^n$ is not yet understood in
closed form; its growth is related to deep questions about Vassiliev invariants.
\end{theorem}

\begin{remark}[The $STU$ relation and Jacobi diagrams]
An important extension replaces chord diagrams by \emph{Jacobi diagrams}
(also called Chinese character diagrams or $3$-graphs): diagrams with
univalent vertices on a circle and trivalent vertices in the interior.
These satisfy the \emph{STU relation}:
\[
  S = T - U,
\]
and the resulting algebra $\mathcal{B}$ is isomorphic to $\Ac$ (the
Poincar\'e-Birkhoff-Witt isomorphism for Jacobi diagrams).  Jacobi diagrams
are more natural for Lie-algebraic weight systems.
\end{remark}

% -------------------------------------------------------
\section{Weight Systems and Lie Algebras}\label{sec:weight}
% -------------------------------------------------------

\begin{definition}[Weight system]
A \emph{weight system of order $n$} is a linear map
$W: \Ac_n \to \mathbf{k}$ (satisfying the $4T$ and $1T$ relations, i.e.\
$W \in \Ac_n^* = \Hom(\Ac_n, \mathbf{k})$).  A weight system of all orders
is a sequence $\{W_n\}$ with $W_n \in \Ac_n^*$ for each $n$, forming an
element of $\Ac^* = \prod_n \Ac_n^*$.
\end{definition}

\begin{theorem}[Vassiliev invariants $\leftrightarrow$ weight systems,
Vassiliev-Kontsevich]\label{thm:vs_ws}
There is a canonical isomorphism of filtered vector spaces:
\[
  \frac{\mathcal{V}_n}{\mathcal{V}_{n-1}} \xrightarrow{\;\sim\;} \Ac_n^*.
\]
Specifically, given a Vassiliev invariant $v \in \mathcal{V}_n$, its
``symbol'' is a well-defined element of $\Ac_n^*$, and every element of
$\Ac_n^*$ arises this way.
\end{theorem}

\begin{proof}[Proof sketch]
The symbol of $v$ is defined by evaluating $v$ on a singular knot with
$n$ double points and recording the resulting chord diagram.  The $4T$
relation arises from isotopy moves on the knot, and the $1T$ relation
corresponds to the framing condition.  Surjectivity follows from the
construction of the Kontsevich integral.
\end{proof}

\begin{theorem}[Lie-algebraic weight systems, Bar-Natan 1995]
\label{thm:lie_ws}
Let $\mathfrak{g}$ be a metrised finite-dimensional Lie algebra (with
invariant non-degenerate bilinear form $(\cdot,\cdot)$), and let $R$ be
a finite-dimensional representation of $\mathfrak{g}$.  Define the
weight system $W_{\mathfrak{g},R}: \Ac_n \to \mathbf{k}$ by:
\[
  W_{\mathfrak{g},R}(D) = \tr_R\!\left(\prod_{\text{chords } (i,j) \text{ of } D}
  \Omega_{ij}\right),
\]
where $\Omega = \sum_a e_a \otimes e^a$ is the Casimir element ($\{e_a\}$
a basis, $\{e^a\}$ the dual basis), and the product is taken in order along
the circle.  Then $W_{\mathfrak{g},R}$ is a well-defined weight system.
\end{theorem}

\begin{example}[Weight systems from $\mathfrak{sl}_2$]
For $\mathfrak{g} = \mathfrak{sl}_2(\C)$ with the 2-dimensional (standard)
representation, the weight system gives the coefficients of the coloured
Jones polynomial $J_2$.  For the $n$-dimensional representation, one obtains
$J_n$ (the $n$-coloured Jones polynomial).
\end{example}

\begin{conjecture}[All weight systems are Lie-algebraic]\label{conj:all_lie}
Every weight system arises from a finite-dimensional metrised Lie algebra
and a representation thereof.  Equivalently, all Vassiliev invariants come
from quantum groups.
\end{conjecture}

This conjecture is widely believed but remains open.

% -------------------------------------------------------
\section{The Kontsevich Integral}\label{sec:KI}
% -------------------------------------------------------

We now construct the Kontsevich integral.  Let $K \subset \R^2 \times \R$
be an oriented knot in ``general position'' with respect to the height
function $t: \R^3 \to \R$ (i.e.\ all critical points of $t|_K$ are
non-degenerate).

\begin{definition}[Kontsevich integral, analytic form]\label{def:KI}
For a knot $K$ in general position with respect to the height function $t$,
set $z = x + iy$ for the first two coordinates.  The \emph{Kontsevich integral}
is:
\[
  Z(K) = \sum_{m=0}^\infty \frac{1}{(2\pi i)^m}
  \int_{t_{\min} < t_1 < \cdots < t_m < t_{\max}}
  \sum_{\substack{P = \{(z_j, z_j')\}_{j=1}^m \\ z_j(t_j), z_j'(t_j) \in K}}
  (-1)^{\downarrow P} D_P \; \bigwedge_{j=1}^m \frac{d z_j - d z_j'}{z_j - z_j'},
\]
where:
\begin{itemize}
  \item The outer sum is over $m = 0, 1, 2, \ldots$
  \item The inner sum is over all ways to choose $m$ horizontal chords
        $(z_j(t_j), z_j'(t_j))$ with distinct heights $t_1 < \cdots < t_m$.
  \item $(-1)^{\downarrow P}$ is $(-1)$ raised to the number of points in
        $P$ that are on a strand of $K$ oriented downward.
  \item $D_P$ is the chord diagram determined by the pairing $P$ (connecting
        the preimages of $z_j$ and $z_j'$ on $S^1$ with chords).
  \item The form $\bigwedge_{j=1}^m \frac{dz_j - dz_j'}{z_j - z_j'}$ is
        integrated over the simplex
        $t_{\min} < t_1 < \cdots < t_m < t_{\max}$.
\end{itemize}
\end{definition}

\begin{theorem}[Kontsevich 1993, invariance]\label{thm:KI_invariant}
The Kontsevich integral $Z(K) \in \hat{\mathcal{A}} = \prod_{n \geq 0} \mathcal{A}_n$
(the completed chord diagram algebra) is an invariant of framed knots.
Modulo the framing correction, it is an invariant of unframed knots.
\end{theorem}

\begin{proof}[Proof sketch]
The proof uses the fact that:
\begin{enumerate}[label=(\arabic*)]
  \item $Z(K)$ is invariant under horizontal isotopies (diffeomorphisms of
        the plane at each height $t$): this uses the fact that
        $\frac{dz - dz'}{z - z'}$ is a closed $1$-form.
  \item $Z(K)$ is invariant under the Reidemeister moves II and III (but
        \emph{not} Reidemeister I, which changes the framing by $1$): this
        requires a careful analysis near critical points, using the associator.
  \item The framing correction uses the isolated chord diagram $\theta$
        (the only chord diagram in $\Ac_1$ before modding out by $1T$).
\end{enumerate}
Full details in Kontsevich~\cite{kontsevich1993} and Bar-Natan~\cite{barnatan1995}.
\end{proof}

\begin{remark}[The Magnus expansion interpretation]
The Kontsevich integral can also be interpreted as the holonomy of a flat
connection on the configuration space $\mathrm{Conf}_n(\C)$ of $n$ points
in the complex plane.  This connection is the \emph{Knizhnik-Zamolodchikov
(KZ) connection}, which takes values in the algebra $\hat{\mathcal{A}}$.
\end{remark}

% -------------------------------------------------------
\section{Convergence and Regularisation}\label{sec:convergence}
% -------------------------------------------------------

The integral in Definition~\ref{def:KI} has logarithmic divergences near
the critical points of $K$.  These must be regularised.

\begin{definition}[Critical point contribution]
Near a local maximum (resp.\ minimum) of $K$ at height $t_0$, the integral
diverges as $\int \frac{dt}{t_0 - t} = -\log(t_0 - t)$.  Regularise by:
\[
  Z(K) = \lim_{\varepsilon \to 0} \left[
    \text{(integral over } t_{\min}+\varepsilon < \cdots < t_{\max}-\varepsilon)
    \cdot \exp\!\left(-\frac{1}{2} w(\varepsilon) \cdot \Phi\right)
  \right],
\]
where $w(\varepsilon)$ counts the ``winding'' near critical points and
$\Phi$ is the Drinfeld associator.
\end{definition}

\begin{theorem}[Le-Murakami 1996, convergence]\label{thm:convergence}
After the regularisation above, the Kontsevich integral converges and
defines a well-defined element of $\hat{\mathcal{A}}$ for any Morse knot
$K$.
\end{theorem}

% -------------------------------------------------------
\section{The Drinfeld Associator $\Phi_{\mathrm{KZ}}$}\label{sec:associator}
% -------------------------------------------------------

The Drinfeld associator is the key algebraic object governing the behaviour
of the Kontsevich integral near crossings and critical points.

\begin{definition}[KZ equation]
The \emph{Knizhnik-Zamolodchikov (KZ) equation} for a function $G: \C\setminus\{0,1\} \to \hat{\mathcal{A}}$ is the differential equation:
\[
  \frac{dG}{dz} = \left(\frac{A}{z} + \frac{B}{z-1}\right) G,
\]
where $A, B$ are formal variables (generators of the algebra $\hat{\mathcal{A}}$
associated to the chord diagrams on two strands).  Here $A$ corresponds to
the chord at $0$ and $B$ to the chord at $1$.
\end{definition}

\begin{theorem}[Drinfeld 1990, Drinfeld associator]\label{thm:drinfeld}
The KZ equation has a unique solution $G_0(z)$ that is asymptotically
$z^A$ as $z \to 0^+$ (on the real axis), and a unique solution $G_1(z)$
that is asymptotically $(1-z)^B$ as $z \to 1^-$.  The \emph{Drinfeld
associator} (or \emph{KZ associator}) is:
\[
  \Phi_{\mathrm{KZ}} = G_0(z)^{-1} G_1(z) \in \hat{\mathcal{A}}.
\]
It is a group-like element in the completed Hopf algebra $\hat{\mathcal{A}}$:
$\Delta(\Phi_{\mathrm{KZ}}) = \Phi_{\mathrm{KZ}} \otimes \Phi_{\mathrm{KZ}}$
(in the appropriate coproduct sense).
\end{theorem}

\begin{proposition}[Properties of $\Phi_{\mathrm{KZ}}$]
\begin{enumerate}[label=(\roman*)]
  \item $\Phi_{\mathrm{KZ}}$ is an associator: it satisfies the pentagon
        equation
        $(\id \otimes \Delta) \Phi \cdot (\Delta \otimes \id) \Phi =
         (\Phi \otimes 1)({\id \otimes \id \otimes \Delta)}\Phi)(1 \otimes \Phi)$
        and the hexagon equations.
  \item $\Phi_{\mathrm{KZ}}$ is group-like: $\Phi_{\mathrm{KZ}} = \exp(\sum_n c_n [\text{Lie element in degree } n])$, where the coefficients $c_n \in \Q$ involve multiple zeta values.
  \item The lowest-degree term of $\log \Phi_{\mathrm{KZ}}$ in degree $3$
        is $\frac{1}{24}[A, B]$ (related to $\zeta(3)$).
  \item Drinfeld showed that any two associators over $\Q$ are conjugate
        by a well-defined element.
\end{enumerate}
\end{proposition}

\begin{remark}[Associator and multiple zeta values]
The coefficients of $\Phi_{\mathrm{KZ}}$ in the basis of Jacobi diagrams
are (regularised) multiple zeta values (MZVs):
$\zeta(n_1, \ldots, n_k) = \sum_{m_1 > \cdots > m_k > 0} m_1^{-n_1} \cdots m_k^{-n_k}$.
This is one of the deep connections between knot theory and number theory.
\end{remark}

% -------------------------------------------------------
\section{The Coloured Jones Polynomial as a Vassiliev Invariant}\label{sec:jones}
% -------------------------------------------------------

\begin{definition}[Coloured Jones polynomial]
For a knot $K$ and an integer $N \geq 1$, the \emph{$N$-coloured Jones
polynomial} $J_N(K; q)$ is the quantum invariant associated to the
$N$-dimensional irreducible representation $V_N$ of the quantum group
$U_q(\mathfrak{sl}_2)$, normalised so that $J_N(\text{unknot}; q) = 1$.
\end{definition}

\begin{theorem}[Vassiliev invariants from the Jones polynomial]
\label{thm:jones_vassiliev}
The $N$-coloured Jones polynomial $J_N(K;e^h)$, expanded as a power series
in $h$ (i.e.\ with $q = e^h$ and $h$ formal):
\[
  J_N(K; e^h) = \sum_{n=0}^\infty v_n^{(N)}(K) h^n,
\]
yields Vassiliev invariants $v_n^{(N)}$ of type $\leq n$ for each $n$.
The weight system of $v_n^{(N)}$ is $W_{\mathfrak{sl}_2, V_N}$
(the weight system from the $N$-dimensional representation of $\mathfrak{sl}_2$).
\end{theorem}

\begin{theorem}[Kontsevich integral specialises to coloured Jones]
\label{thm:KI_jones}
Let $W_{\mathfrak{sl}_2, V_N}: \hat{\mathcal{A}} \to \C[[h]]$ be the
Lie-algebraic weight system.  Then:
\[
  W_{\mathfrak{sl}_2, V_N}(Z(K)) = J_N(K; e^h).
\]
That is, the coloured Jones polynomial is the specialisation of the
Kontsevich integral under the weight system $W_{\mathfrak{sl}_2, V_N}$.
\end{theorem}

\begin{remark}[Volume conjecture and the coloured Jones polynomial]
The \emph{volume conjecture} (Kashaev 1995, Murakami-Murakami 2001) states
that for a hyperbolic knot $K$:
\[
  \lim_{N \to \infty} \frac{2\pi}{N} \log|J_N(K; e^{2\pi i/N})| = \mathrm{vol}(S^3 \setminus K),
\]
where $\mathrm{vol}$ denotes the hyperbolic volume.  This connects the
Vassiliev/Kontsevich world to hyperbolic geometry.
\end{remark}

% -------------------------------------------------------
\section{Universality of the Kontsevich Integral}\label{sec:universality}
% -------------------------------------------------------

\begin{theorem}[Kontsevich, universality]\label{thm:universality}
The Kontsevich integral $Z: \{\text{knots}\} \to \hat{\mathcal{A}}$ is the
\emph{universal Vassiliev invariant}: every Vassiliev invariant factors
through $Z$.  More precisely, for every weight system $W \in \hat{\mathcal{A}}^*$,
the composite $W \circ Z: \{\text{knots}\} \to \mathbf{k}$ is a Vassiliev
invariant, and every Vassiliev invariant arises in this way.
\end{theorem}

\begin{proof}
By Theorem~\ref{thm:vs_ws}, Vassiliev invariants are classified by weight
systems.  Given a weight system $W$, define $v_W = W \circ Z$.  Then
$v_W$ is a Vassiliev invariant: its degree-$n$ part is $W_n \circ Z_n$,
where $Z_n$ is the degree-$n$ part of $Z$.  Conversely, any $v \in \mathcal{V}_n$
has a symbol $\bar{v} \in \Ac_n^*$, and one shows that $v = \bar{v} \circ Z_n$
modulo lower-order terms (by induction on $n$, using that $Z_n$ spans $\Ac_n$
over rational functions of $K$).
\end{proof}

\begin{corollary}
The Kontsevich integral encodes \emph{all} quantum group invariants of knots
(all coloured Jones polynomials, HOMFLY-PT polynomial, Kauffman polynomial, etc.)
as specialisations.
\end{corollary}

% -------------------------------------------------------
\section{The Completeness Conjecture}\label{sec:completeness}
% -------------------------------------------------------

\begin{conjecture}[Kontsevich integral separates knots]\label{conj:complete}
The Kontsevich integral $Z: \{\text{knots}\} / \text{isotopy} \to \hat{\mathcal{A}}$
is injective.  That is:
\[
  Z(K_1) = Z(K_2) \implies K_1 \text{ and } K_2 \text{ are isotopic.}
\]
\end{conjecture}

\begin{remark}
This would mean that Vassiliev invariants form a \emph{complete} system of
knot invariants: they collectively distinguish all knots.  The conjecture
is widely believed but completely open.  Even the weaker statement that
Vassiliev invariants distinguish all knots from the unknot is unknown.
\end{remark}

\begin{proposition}[Evidence for Conjecture~\ref{conj:complete}]
\begin{enumerate}[label=(\roman*)]
  \item The Kontsevich integral distinguishes all \emph{prime} knots up to
        at least 10 crossings (verified computationally by Bar-Natan~\cite{barnatan1995}).
  \item $Z(K)$ distinguishes knots from their mirror images if $K$ is
        chiral (e.g.\ the trefoil).
  \item No two distinct knots with fewer than $10$ crossings have been found
        with the same Kontsevich integral.
  \item Partial converse: if $Z(K) = Z(\text{unknot})$ and $K$ is a knot
        in a homology sphere, then $K$ has all Vassiliev invariants equal
        to those of the unknot.
\end{enumerate}
\end{proposition}

\begin{conjecture}[Vassiliev invariants detect the unknot]\label{conj:unknot}
If all Vassiliev invariants of $K$ agree with those of the unknot, then
$K$ is the unknot.
\end{conjecture}

This is a weaker form of Conjecture~\ref{conj:complete} and is also open.

% -------------------------------------------------------
\section{The Wheeling Theorem}\label{sec:barley}
% -------------------------------------------------------

The Wheeling theorem, proved by Bar-Natan, Le, and Thurston~\cite{barnatan2003},
is one of the deepest results in the theory of Vassiliev invariants.

\begin{definition}[The wheel $\omega$ and the wheeling map]
In the algebra $\mathcal{B}$ of Jacobi diagrams, the \emph{wheel} $\omega_n$
is the diagram consisting of a ``hub'' with $n$ spokes.  Specifically,
$\omega_n$ has $n$ univalent vertices connected to a single trivalent vertex
(the hub) by $n$ edges.  Define the element
\[
  \Omega = \sum_{n=0}^\infty \frac{b_{2n}}{n!} \omega_{2n}
  \in \hat{\mathcal{B}},
\]
where $b_{2n}$ are the modified Bernoulli numbers ($\frac{t}{e^t-1} = \sum_n b_n t^n/n!$).
The \emph{wheeling map} $\Phi_\Omega: \mathcal{B} \to \mathcal{A}$ is defined
by inserting the wheels $\Omega$ at appropriate legs.
\end{definition}

\begin{theorem}[Wheeling theorem, Bar-Natan--Le--Thurston 2003]\label{thm:wheeling}
The wheeling map
\[
  \Phi_\Omega: (\mathcal{B}, \cdot_{\text{disjoint}}) \to (\mathcal{A}, \cdot_{\text{connected sum}})
\]
is an isomorphism of graded algebras.  Here $\mathcal{B}$ carries the
disjoint union product and $\mathcal{A}$ carries the connected sum product.
\end{theorem}

\begin{proof}[Proof idea]
The proof proceeds by showing that $\Phi_\Omega$ is an algebra map using
the combinatorics of the STU relation and the properties of Bernoulli numbers.
The key step is an explicit formula relating the two products via a ``wheeling''
operation that inserts the element $\Omega$ to correct for the difference
between disjoint and connected-sum products.
The full proof uses the language of props and operads, and occupies
\cite{barnatan2003}.
\end{proof}

\begin{corollary}[Duflo isomorphism as a specialisation]
The Wheeling theorem is a ``categorification'' of the Duflo isomorphism in Lie
theory.  Specifically, applying the weight system $W_{\mathfrak{g}}$ for a
Lie algebra $\mathfrak{g}$ to the Wheeling theorem recovers the Duflo
isomorphism $\mathrm{Sym}(\mathfrak{g})^{\mathfrak{g}} \cong Z(U\mathfrak{g})$.
\end{corollary}

% -------------------------------------------------------
\section{The Le-Murakami-Ohtsuki Invariant}\label{sec:LMO}
% -------------------------------------------------------

The Le-Murakami-Ohtsuki (LMO) invariant generalises the Kontsevich integral
to $3$-manifolds.

\begin{definition}[Surgery presentation and the LMO invariant]
Any closed orientable $3$-manifold $M$ can be obtained by surgery on a link
$L \subset S^3$.  The \emph{LMO invariant} $Z_{\mathrm{LMO}}(M) \in \hat{\mathcal{A}}$
is defined by:
\[
  Z_{\mathrm{LMO}}(M) = \iota_n\left(Z(L) / \iota_n(Z(\bigcirc))^{\otimes n}\right)^{\otimes \sigma^\pm},
\]
where $n$ is the number of components of $L$, $\sigma^\pm$ is the signature
of the linking matrix, $\bigcirc$ is the unknot, and $\iota_n$ is a specific
map that ``deframes'' the invariant.  See Le-Murakami-Ohtsuki~\cite{lmo1998}
for the precise definition.
\end{definition}

\begin{theorem}[Le-Murakami-Ohtsuki 1998]\label{thm:LMO}
The LMO invariant $Z_{\mathrm{LMO}}(M)$ is a well-defined invariant of
closed oriented $3$-manifolds.  It dominates all quantum invariants of
$3$-manifolds in the sense that the Reshetikhin-Turaev invariants (for all
quantum groups) are specialisations of $Z_{\mathrm{LMO}}$.
\end{theorem}

\begin{remark}
The LMO invariant is conjectured to be a complete invariant of rational
homology spheres (up to orientation-preserving diffeomorphism), but this
is open.
\end{remark}

% -------------------------------------------------------
\section{Relation to Chern-Simons Theory}\label{sec:chern}
% -------------------------------------------------------

Witten~\cite{witten1989} proposed that all quantum invariants of knots and
$3$-manifolds arise from the path integral:
\[
  Z_{\mathrm{CS}}(M, K) = \int_{\mathcal{A}} e^{\frac{ik}{4\pi} \int_M \tr\!(A \wedge dA + \tfrac{2}{3} A \wedge A \wedge A)} W_R(K) \,\mathcal{D}A,
\]
where $A$ is a gauge field (connection), $W_R(K) = \tr_R \mathrm{Hol}_A(K)$
is the Wilson loop observable in representation $R$, and the integral is over
all connections $\mathcal{A}$ modulo gauge equivalence.

\begin{theorem}[Axiomatic reformulation, Reshetikhin-Turaev 1991]
The Chern-Simons path integral, made rigorous via the Reshetikhin-Turaev
construction~\cite{reshetikhin1991} using quantum groups at roots of unity,
gives topological invariants of $3$-manifolds and knots therein.
\end{theorem}

\begin{proposition}[Kontsevich integral as perturbative Chern-Simons]
The Kontsevich integral $Z(K)$ arises as the perturbative expansion of the
Chern-Simons path integral around the trivial connection:
\[
  Z(K) = \sum_{n=0}^\infty \hbar^n \cdot [\text{Feynman diagram contribution of order } n].
\]
The Feynman diagrams are precisely Jacobi diagrams (the $3$-valent graphs
correspond to $\tr([A,A] \cdot A)$ vertices, and propagators give the
$1$-forms $\frac{dz-dz'}{z-z'}$).
\end{proposition}

\begin{remark}[Physical interpretation]
The $4T$ relation in chord diagram algebra corresponds to the Jacobi identity
for the Lie bracket in $\mathfrak{g} = \mathrm{Lie}(G)$, arising from the
antisymmetry of the structure constants $f_{abc}$.  This is why Lie-algebraic
weight systems are natural: the gauge field $A$ takes values in $\mathfrak{g}$.
\end{remark}

% -------------------------------------------------------
\section{Bar-Natan's Computations and Applications}\label{sec:barnatan}
% -------------------------------------------------------

\begin{theorem}[Bar-Natan 1995, generators of $\mathcal{A}$]
\label{thm:generators}
The chord diagram algebra $\mathcal{A}$ is a polynomial algebra:
\[
  \mathcal{A} \cong \Q[w_2, w_3, w_4, \ldots],
\]
where $w_n \in \mathcal{A}_n$ are the \emph{wheel elements}.  More precisely,
$\mathcal{A}$ is a graded polynomial algebra on generators indexed by the
wheel diagrams $w_n$ for $n \geq 2$.
\end{theorem}

\begin{remark}
This is a theorem about $\mathcal{A}$ as an \emph{algebra} (with connected-sum
product).  It implies that the graded Euler characteristic of $\mathcal{A}$
is $\prod_{n \geq 2} (1 - t^n)^{-1}$, which equals the generating function
for the number of partitions with parts $\geq 2$.
\end{remark}

\begin{example}[Explicit computation of $Z$ for the trefoil]
For the left-hand trefoil knot $3_1$, the Kontsevich integral begins:
\[
  Z(3_1) = 1 + a_2 w_2 + a_3 w_3 + (a_{22} w_2^2 + a_4 w_4) + \cdots,
\]
where $a_2 = \frac{1}{2}$, $a_3 = \frac{1}{12}$, and higher coefficients
involve multiple zeta values.  The coefficient $a_2$ recovers the Casson
invariant of the trefoil.
\end{example}

% -------------------------------------------------------
\section{Open Problems}\label{sec:open}
% -------------------------------------------------------

We collect the major open problems in this area:

\begin{conjecture}[Completeness of Vassiliev invariants, Birman-Lin 1993]
\label{conj:complete2}
Vassiliev invariants collectively separate all knots.  That is, if two
knots $K_1, K_2$ satisfy $v(K_1) = v(K_2)$ for all Vassiliev invariants
$v$, then $K_1$ and $K_2$ are isotopic.
\end{conjecture}

\begin{conjecture}[Finiteness of the dimension of $\mathcal{A}_n$]
\label{conj:finiteness}
Although $\dim \mathcal{A}_n < \infty$ is known for each $n$, the precise
growth of $\dim \mathcal{A}_n$ as $n \to \infty$ is unknown.  The conjecture
is that $\sum_n \dim(\mathcal{A}_n) t^n$ has a natural boundary on the unit
circle.
\end{conjecture}

\begin{conjecture}[All weight systems are Lie-algebraic, cf.\ Conjecture~\ref{conj:all_lie}]
Every weight system arises from a metrised Lie algebra and a representation.
This would imply that all Vassiliev invariants come from quantum groups.
\end{conjecture}

\begin{conjecture}[Volume conjecture, Kashaev-Murakami-Murakami]\label{conj:volume}
For a hyperbolic knot $K$:
\[
  \lim_{N \to \infty} \frac{2\pi}{N} \log |J_N(K; e^{2\pi i/N})| = \mathrm{vol}(S^3 \setminus K).
\]
\end{conjecture}

\begin{remark}[Relation to categorification]
Khovanov homology~\cite{khovanov2000} is a categorification of the Jones
polynomial: it provides a bigraded homology theory $\mathrm{Kh}^{i,j}(K)$
such that $J_2(K;q) = \sum_{i,j} (-1)^i q^j \dim \mathrm{Kh}^{i,j}(K)$.
There is no known categorification of the full Kontsevich integral,
though partial categorifications (e.g.\ Khovanov-Rozansky homology) exist.
This is a major open problem.
\end{remark}

\begin{thebibliography}{99}

\bibitem{vassiliev1990}
V.\ A.\ Vassiliev,
\textit{Cohomology of knot spaces},
in: Theory of Singularities and Its Applications (V.\ I.\ Arnold, ed.),
Adv.\ Soviet Math.\ \textbf{1}, Amer.\ Math.\ Soc., 1990, pp.\ 23--69.

\bibitem{kontsevich1993}
M.\ Kontsevich,
\textit{Vassiliev's knot invariants},
in: I.\ M.\ Gel'fand Seminar, Adv.\ Soviet Math.\ \textbf{16}, Part 2,
Amer.\ Math.\ Soc., Providence, RI, 1993, pp.\ 137--150.

\bibitem{barnatan1995}
D.\ Bar-Natan,
\textit{On the Vassiliev knot invariants},
Topology \textbf{34} (1995), 423--472.

\bibitem{jones1985}
V.\ F.\ R.\ Jones,
\textit{A polynomial invariant for knots via von Neumann algebras},
Bull.\ Amer.\ Math.\ Soc.\ \textbf{12} (1985), 103--111.

\bibitem{witten1989}
E.\ Witten,
\textit{Quantum field theory and the Jones polynomial},
Comm.\ Math.\ Phys.\ \textbf{121} (1989), 351--399.

\bibitem{drinfeld1990}
V.\ G.\ Drinfeld,
\textit{Quasi-Hopf algebras},
Leningrad Math.\ J.\ \textbf{1} (1990), 1419--1457.

\bibitem{birmanlin1993}
J.\ Birman, X.\ Lin,
\textit{Knot polynomials and Vassiliev's invariants},
Invent.\ Math.\ \textbf{111} (1993), 225--270.

\bibitem{barnatan2003}
D.\ Bar-Natan, S.\ Le, D.\ Thurston,
\textit{Two applications of elementary knot theory to Lie algebras and Vassiliev invariants},
Geom.\ Topol.\ \textbf{7} (2003), 1--31.

\bibitem{lmo1998}
T.\ Le, J.\ Murakami, T.\ Ohtsuki,
\textit{On a universal perturbative invariant of $3$-manifolds},
Topology \textbf{37} (1998), 539--574.

\bibitem{reshetikhin1991}
N.\ Reshetikhin, V.\ G.\ Turaev,
\textit{Invariants of $3$-manifolds via link polynomials and quantum groups},
Invent.\ Math.\ \textbf{103} (1991), 547--597.

\bibitem{khovanov2000}
M.\ Khovanov,
\textit{A categorification of the Jones polynomial},
Duke Math.\ J.\ \textbf{101} (2000), 359--426.

\bibitem{leturc1997}
T.\ Q.\ Le, J.\ Murakami,
\textit{The universal Vassiliev-Kontsevich invariant for framed oriented links},
Compositio Math.\ \textbf{102} (1996), 41--64.

\bibitem{barnatan1998}
D.\ Bar-Natan,
\textit{Non-associative tangles},
in: Geometric Topology (Proc.\ 1993 Georgia Internat.\ Topology Conf.),
AMS/IP Studies in Adv.\ Math.\ \textbf{2}, Amer.\ Math.\ Soc., 1997, pp.\ 139--183.

\bibitem{chmutov2012}
S.\ Chmutov, S.\ Duzhin, J.\ Mostovoy,
\textit{Introduction to Vassiliev Knot Invariants},
Cambridge University Press, Cambridge, 2012.

\bibitem{ohtsuki2002}
T.\ Ohtsuki,
\textit{Quantum Invariants: A Study of Knots, 3-Manifolds, and Their Sets},
Series on Knots and Everything \textbf{29}, World Scientific, 2002.

\bibitem{murakami2001}
H.\ Murakami, J.\ Murakami,
\textit{The colored Jones polynomials and the simplicial volume of a knot},
Acta Math.\ \textbf{186} (2001), 85--104.

\end{thebibliography}

\end{document}
